\paragraph{Problem 3 answer}\footnote
{\color{red} MAKE SURE TO READ THE COMMENTS IN problem\_3\_answer.tex}
\GradescopeComment{This should be on page 5 (top) of your PDF.}

\begin{enumerate}

\item[(a)] Here is a definition of all the subproblems the algorithm will solve:

  \begin{blue}

    \REPLACEME 
    
    % Here's the corresponding answer for Making Change.
    % Model your answer on it.

    % For $t\in\{0,1,\ldots, T\}$, define $M(t)$ to be the minimum number of coins
    % needed to make change for $t$.  The final answer is $M(T)$.

  \end{blue}
  
\item[(b)] Here's the recurrence relation:

  \begin{blue}

    \REPLACEME 

    % Here's the corresponding answer for Making Change.
    % Model your answer on it.

    \[
      % M(t) =
      % \begin{cases}
      %   0
      %   & \text{if } t=0
      %   \\
      %   \min\{ 1 + M(t-d_i) : d_i \in D,\, d_i \le t\}
      %   & \text{otherwise.}
      % \end{cases}
    \]

  \end{blue}
\end{enumerate}  

\begin{itemize}
\item[(c)]  Here is pseudo-code for an iterative implementation (don't forget the return statement):
\end{itemize}

\begin{blue}
  
  \REPLACEME 

  % Here's the corresponding answer for Making Change.
  % Model your answer on it.

  % \begin{myframed}
  %   \begin{steps}

  %   \item[] \textsf{makeChange}$(D=\{d_1, d_2, \ldots, d_k\}, T)$:
      
  %     \step For $t=0,1,\ldots, T$:

  %     \step~~~\(
  %     m[t] =
  %     \begin{cases}
  %       0
  %       & \text{if } t=0
  %       \\
  %       \min\{ 1 + m[t-d_i] : d_i \in D,\, d_i \le t\}
  %       & \text{otherwise.}
  %     \end{cases}
  %     \)

  %     \step return $m[T]$
      
  %   \end{steps}
  % \end{myframed}

\end{blue}


\vfill

\paragraph{Problem \arabic{problem} collaboration and citations}~\GradescopeComment{This should be on page 7 (bottom) of your PDF.}

I collaborated on this problem with the following people:
\begin{blue}
  \begin{itemize}
  \item \REPLACEME                    % write "none" if none
  \end{itemize}
\end{blue}

My answers use ideas from the following sources (other than the lecture notes and textbooks): 
\begin{blue}
  \begin{itemize}
  \item \REPLACEME                    % write "none" if none
  \end{itemize}
\end{blue}



%%% Local Variables:
%%% mode: latex
%%% TeX-master: "homework_7"
%%% End:
